\section{第1章 函数与极限}
\subsection{函数}
\begin{exercise}{1}{定义域与参数约束}
求 $f(x)=\arccos(2x+1)+\log_2(3x+2)$ 的定义域。再确定实数 $a$ 的取值范围，使区间 $[a,a+\frac14]$ 全部包含于该定义域。
\end{exercise}
\begin{exercise}{2}{复合函数的定义域}
设 $f(x)=\frac{x-1}{x+1}$。求 $f(-2)$、$f(f(x))$，并写出复合函数的完整定义域；判断化简后的表达式能否在其自然定义域上与原复合函数视为同一函数。
\end{exercise}
\begin{exercise}{3}{换元与函数确定范围}
已知对一切 $x\ne0$，$f(1/x)=4/x+2x^2+1$。求 $f(t)$ 的已知表达式及其适用范围，再求 $f(x^2+1)$；说明上述条件是否确定了 $f(0)$。
\end{exercise}
\begin{exercise}{4}{反三角函数与值域}
设 $f(x)=\arcsin\frac4x$。求其定义域及 $f(-4)$、$f(8/\sqrt2)$、$f(8/\sqrt3)$，并确定其值域及各端点是否能够取到。
\end{exercise}
\begin{exercise}{5}{复合函数值域}
求 $f(x)=\sin\frac{\pi x}{2(1+x^2)}$ 的值域及取到最大值、最小值的全部自变量，并说明为何不能直接把正弦函数的值域作为答案。
\end{exercise}
\begin{exercise}{6}{反三角函数的主值}
设 $f(x)=\arccos2x$、$g(x)=\arctan x$、$h(x)=\operatorname{arccot}(x/2)$，其中 $\operatorname{arccot}$ 取值于 $(0,\pi)$。求 $f(\sqrt3/4)g(-1)h(2)$；分别写出三个函数的定义域和值域。
\end{exercise}
\begin{exercise}{7}{非单射换元与定义域}
已知对一切非零实数 $x$，$f(x+1/x)=\frac{4x^2}{x^4+1}$。求条件所能确定的 $f(t)$ 及 $t$ 的范围，并验证 $x$ 与 $1/x$ 给出的函数值一致；说明能否据此确定 $f(0)$。
\end{exercise}
\begin{exercise}{8}{分段函数与奇偶性}
设 $f(x)=\begin{cases}\frac{x}{\pi}|\sin x|,&|x|<\pi/3,\\0,&|x|\ge\pi/3.\end{cases}$ 求 $f(\pi/6)$、$f(-\pi/4)$、$f(2)$；判断其奇偶性，并讨论 $x=\pm\pi/3$ 两点处能否通过重新赋值使函数连续。
\end{exercise}
\begin{exercise}{9}{反解复合函数的分支限制}
设 $f(t)=\tan2t$，$f(g(x))=x^2-1$，且 $|g(x)|\le\pi/8$。求 $g(x)$ 及其最大定义域，说明范围限制在确定反函数分支和定义域时各起什么作用。
\end{exercise}
\begin{exercise}{10}{最大值函数的分段表达}
设 $\varphi(x)=\max\{x,x^2,1\}$。求定义域、分段表达式以及 $\varphi(-2)$、$\varphi(-1)$、$\varphi(0)$；确定该函数的最小值与全部取值点。
\end{exercise}
\begin{exercise}{11}{复合次序与分段边界}
设 $f(x)=\begin{cases}-1,&x<0,\\1,&x\ge0,\end{cases}$、$g(x)=2x-1$。分别写出 $f(g(x))$、$g(f(x))$ 的分段表达式，并求两者相等的全部 $x$。
\end{exercise}
\begin{exercise}{12}{复合映射与值域约束}
设 $f(x)=\begin{cases}2-x,&x<0,\\2+x,&x\ge0,\end{cases}$、$g(x)=\begin{cases}x,&x<0,\\x^2,&x\ge0.\end{cases}$ 求 $f(g(x))$、$g(f(x))$ 的分段表达式和值域，并判断它们是否为偶函数。
\end{exercise}
\begin{exercise}{13}{分段函数的四则运算}
设 $f(x)=\begin{cases}2x+1,&x<0,\\x^2+1,&x\ge0,\end{cases}$、$g(x)=\begin{cases}\sin x,&x<\pi,\\\cos x,&x\ge\pi.\end{cases}$ 写出 $f+g$、$fg$ 的分段表达式，并分别讨论二者在 $0$ 与 $\pi$ 处的连续性。
\end{exercise}
\begin{exercise}{14}{反函数及端点范围}
设 $f(x)=\frac{1-\sqrt{1+x}}{1+\sqrt{1+x}}$。确定定义域和值域，证明该函数可逆并求反函数；写出反函数的定义域，解释原函数值域中未包含的端点。
\end{exercise}
\begin{exercise}{15}{奇偶性的定义域条件}
设 $a>0$。分别判断下列函数的奇偶性，写明定义域，并对非奇非偶的函数给出依据：$ (x^2+1)\sin^3x$，$\ln(x+\sqrt{x^2+a^2})$，$\frac{2^{-x}-2^x}{2^{-x}+2^x}$，$x^2\sec x+\csc^2x$。讨论第二个函数在何种 $a$ 下成为奇函数。
\end{exercise}
\begin{exercise}{16}{抽象函数复合的奇偶性}
设 $f,g:\mathbb R\to\mathbb R$，$f$ 为奇函数，$g$ 为偶函数。判断 $f\circ f$、$f\circ g$、$g\circ f$、$g\circ g$ 必然具有的奇偶性并证明。说明“偶函数一定不是奇函数”是否成立。
\end{exercise}
\begin{exercise}{17}{单调性的直接证明}
不用导数，证明 $f(x)=x^3+2x+3$ 在 $\mathbb R$ 上严格单调；进一步确定参数 $a$ 的全部取值，使 $x^3+ax+3$ 仍在 $\mathbb R$ 上严格递增。
\end{exercise}
\begin{exercise}{18}{区间上的有界性}
分别判断 $f(x)=\frac1{x-2}$ 在 $(2,3)$ 和 $(3,+\infty)$ 上是否有界；有界时求上确界与下确界，无界时按定义构造使函数值任意大的点。
\end{exercise}
\begin{exercise}{19}{最小正周期}
求 $f(x)=2+\cos^2(2x-1)$ 的最小正周期，并证明不存在更小的正周期。讨论常数项与相位平移是否影响该结论。
\end{exercise}
\begin{exercise}{20}{取整函数与周期}
记 $\lfloor x\rfloor$ 为不超过 $x$ 的最大整数。对 $f(x)=x-\lfloor x\rfloor$，确定其值域、最小正周期，并判断有界性、单调性和奇偶性；写出 $f(-x)$ 与 $f(x)$ 在非整数点的关系。
\end{exercise}
\begin{exercise}{21}{函数关系与反函数}
设 $f:(0,+\infty)\to\mathbb R$ 满足 $f(x)+2f(1/x)=x$。求 $f(x)$；选取包含 $1$ 的最大区间，使 $f$ 在该区间上一一对应，并求其反函数及定义域。
\end{exercise}
\begin{exercise}{22}{周期函数的有界条件}
设 $f:\mathbb R\to\mathbb R$ 以 $T>0$ 为周期。证明 $f$ 在 $\mathbb R$ 上有界当且仅当它在 $[0,T)$ 上有界；构造一个无界周期函数，并说明“周期函数必有界”遗漏了什么条件。
\end{exercise}
\subsection{数列的极限}
\begin{exercise}{23}{数列极限定义}
用 $\varepsilon$--$N$ 定义证明 $\lim_{n\to\infty}\frac{n}{2n+1}=\frac12$，给出可取的 $N(\varepsilon)$；再求误差小于 $10^{-3}$ 时所需的最小正整数 $n$。
\end{exercise}
\begin{exercise}{24}{幂数列的参数分类}
对全部实数 $x$，讨论 $x^n$ 的有限极限、正负无穷趋势或振荡情形。对有限极限不存在的情形，用子列或绝对值增长给出理由。
\end{exercise}
\begin{exercise}{25}{实指数数列}
对全部实数 $x$，讨论 $n^x$ 在 $n\to\infty$ 时的变化。对 $x<0$ 的情形，用极限定义给出 $N(\varepsilon,x)$；说明固定 $x$ 的分类结论能否直接用于 $x_n=-1/\ln(n+1)$。
\end{exercise}
\begin{exercise}{26}{交错符号数列的定义证明}
设 $x_{2k-1}=-1/(2k)$、$x_{2k}=2^{-k}$（$k\ge1$）。判断数列的极限是否存在并用定义证明；比较两种符号的子列各自趋零的速度，说明交错符号与不同衰减速度是否阻碍收敛。
\end{exercise}
\begin{exercise}{27}{奇偶子列与收敛}
设 $x_n=\begin{cases}(n+1)/(n+2),&n\text{ 为偶数},\\1,&n\text{ 为奇数}.\end{cases}$ 求其极限，并用统一的误差估计验证；若将奇数项改为常数 $a$，求数列收敛的充要条件。
\end{exercise}
\begin{exercise}{28}{交错排列的子列判据}
设 $x_n\to-1$、$y_n\to1$，把它们交错排列成 $z_{2n-1}=x_n$、$z_{2n}=y_n$。判断 $z_n$ 是否收敛，并求它的全部子列极限；进一步证明一般情形 $x_n\to a,y_n\to b$ 下，$z_n$ 收敛当且仅当 $a=b$。
\end{exercise}
\begin{exercise}{29}{收敛性质的条件辨析}
判断并证明或举反例：$x_n\to a$ 则 $|x_n|\to|a|$；$|x_n|\to|a|$ 则 $x_n\to a$；$x_n\to a$ 则 $x_{n+1}\to a$；$x_n\to a$ 则 $1/x_n\to1/a$。对最后一个断言补充使其成立的必要限制。
\end{exercise}
\begin{exercise}{30}{周期振荡数列的子列}
证明 $x_n=(1+1/n)\sin(n\pi/4)$ 不收敛；求其全部子列极限，并说明因子 $1+1/n\to1$ 为何不能消除周期振荡。
\end{exercise}
\begin{exercise}{31}{收敛与连续变换}
设实数列 $x_n$ 满足 $x_n^2\to1$。证明 $|x_n|\to1$；给出 $x_n$ 不收敛的例子，并证明增加条件“$x_n$ 从某项起同号”后，$x_n$ 必收敛。
\end{exercise}
\begin{exercise}{32}{差分趋零的适用范围}
若 $x_n\to a$，证明 $x_{n+1}-x_n\to0$。分别构造无界数列和有界数列，说明该结论的逆命题均不成立；有界例子须给出两个极限不同的子列。
\end{exercise}
\subsection{函数的极限}
\begin{exercise}{33}{函数极限的定义证明}
用 $\varepsilon$--$\delta$ 定义证明 $\lim_{x\to1/2}\frac{1-4x^2}{2x-1}=-2$，给出明确的 $\delta(\varepsilon)$。若在 $x=1/2$ 任意补定义一个值，说明该极限是否改变。
\end{exercise}
\begin{exercise}{34}{无穷远左右方向}
分别求 $e^{2x-1}$ 与 $\operatorname{arccot}(1-x)$ 在 $x\to-\infty$、$x\to+\infty$ 时的极限或无穷趋势，其中 $\operatorname{arccot}\in(0,\pi)$；判断各函数是否存在 $x\to\infty$（即 $|x|\to\infty$）时的有限极限。
\end{exercise}
\begin{exercise}{35}{分段点的左右极限}
设 $f(x)=\begin{cases}(x+1)\sin x,&x<0,\\1-\cos x,&x\ge0.\end{cases}$ 求 $x=0$ 处的左右极限及双侧极限；若把左支改为 $(x+1)\sin x+a$，确定双侧极限存在的全部 $a$。
\end{exercise}
\begin{exercise}{36}{分段极限与定义域方向}
设 $f(x)=\begin{cases}\ln(1-x),&x<1,\\1/\sqrt{2x+1},&x\ge1.\end{cases}$ 求 $x\to1$、$x\to0$、$x\to2$ 及 $|x|\to\infty$ 时的极限或变化趋势，讨论时须区分各方向。
\end{exercise}
\begin{exercise}{37}{跳跃、可去点与无穷远极限}
设 $f(x)=\begin{cases}x+1,&x<0,\\3,&x=0,\\2-x,&0<x<1,\\0,&x=1,\\1+1/x,&x>1.\end{cases}$ 求 $x\to0^-$、$x\to0^+$、$x\to1$、$x\to2$ 及 $x\to+\infty$ 时的极限，并指出哪些结论不能用直接代入函数值获得。
\end{exercise}
\begin{exercise}{38}{局部保号性与有界性}
设 $\lim_{x\to1}f(x)/x=-1$。证明存在 $1$ 的去心邻域使 $f(x)<0$ 且有界；举例说明“任意去心邻域内均处处 $f(x)<0$”与“任意去心邻域内均有界”不由条件保证。
\end{exercise}
\begin{exercise}{39}{极限不等式与量词}
设 $f(x)\to A$、$g(x)\to B$（$x\to x_0$）。证明：若某去心邻域内 $f(x)<g(x)$，则 $A\le B$；若 $A<B$，则某去心邻域内 $f(x)<g(x)$。分别举例说明第一结论不能加强为 $A<B$，第二结论不能扩大到所有 $x$；说明极限存在是否要求 $f(x_0)$ 有定义。
\end{exercise}
\begin{exercise}{40}{无穷远振荡与海涅定理}
用两列趋于 $+\infty$ 的点证明 $x\cos x$ 没有有限极限，并证明它也不趋于绝对值无穷大；再判断 $\cos x/x$ 在同一变化过程中的极限。
\end{exercise}
\begin{exercise}{41}{局部有界性的准确范围}
用极限定义证明：若 $\lim_{x\to x_0}f(x)=A\in\mathbb R$，则 $f$ 在 $x_0$ 的某去心邻域内有界。若 $f(x_0)$ 也有定义，证明可把去心邻域改为邻域；举例说明局部有界不保证极限存在。
\end{exercise}
\begin{exercise}{42}{复合极限的条件辨析}
设 $\lim_{u\to0}f(u)=2$、$\lim_{x\to0}g(x)=0$。构造例子说明未必有 $f(g(x))\to2$；证明若 $g(x)$ 在 $0$ 的某去心邻域内处处不等于 $0$，该复合极限成立，并讨论把这一条件换成 $f(0)=2$ 的情形。
\end{exercise}
\begin{exercise}{43}{海涅定理与局部振荡}
定义 $f(x)=\sin(1/x)+ax$（$x\ne0$），其中 $a\in\mathbb R$。用两列趋于 $0$ 的非零点证明 $f$ 在 $0$ 处无极限；求 $xf(x)$ 在 $0$ 处的极限，并给出相应的 $\varepsilon$--$\delta$ 估计。
\end{exercise}
\subsection{无穷小与无穷大}
\begin{exercise}{44}{有界振荡与无穷小}
证明 $2x\sin(1/x)$ 在 $x\to0$ 时为无穷小。对参数 $p\in\mathbb R$，确定 $|x|^p\sin(1/x)$ 为无穷小的充要条件，并对不成立的参数给出趋零点列。
\end{exercise}
\begin{exercise}{45}{无界与无穷大的区别}
判断 $x\cos x$ 在 $\mathbb R$ 上是否无界，以及在 $x\to+\infty$ 时是否为无穷大（按绝对值趋于无穷解释）；分别构造点列证明，并说明这两个概念的量词差异。
\end{exercise}
\begin{exercise}{46}{变化过程与单侧趋势}
对 $f(x)=e^{1/x}$ 与 $g(x)=x\arctan x$，分别讨论 $x\to0^-$、$x\to0^+$、$x\to-\infty$、$x\to+\infty$ 时的极限或无穷趋势，指出哪些过程中它们是无穷小或无穷大。
\end{exercise}
\begin{exercise}{47}{无穷大与局部有界}
设 $f(x)=\frac{2x}{x-2}$、$g(x)=\ln(x+1)$。分别讨论 $x\to0$、$x\to2$、$x\to-1^+$、$x\to+\infty$ 时两函数的无穷小、无穷大或局部有界性质；对 $x\to2$ 须区分有符号极限与绝对值趋于无穷。
\end{exercise}
\begin{exercise}{48}{乘积趋零的充分条件}
设 $x_ny_n\to0$。判断下列断言并证明或举反例：$x_n$ 发散则 $y_n$ 收敛；$x_n$ 无界则 $y_n$ 有界；$x_n$ 有界则 $y_n\to0$；若 $x_n$ 最终非零且 $1/x_n\to0$，则 $y_n\to0$。
\end{exercise}
\begin{exercise}{49}{有界因子估计}
求 $\lim_{x\to0}\sin x(2+\cos(3/x))$ 与 $\lim_{|x|\to\infty}\frac{\arctan|2x+1|}{2x+1}$，分别写出能保证结论的绝对值估计；讨论第一个式子中的 $2$ 换成任意固定实数后结论是否改变。
\end{exercise}
\begin{exercise}{50}{无穷小的判定}
判定下列各量是否为无穷小，并给出相应极限或增长依据：$\frac{x^2}{\sqrt{x^3+1}}$（$x\to+\infty$），$\frac{x^3-1}{2x^4+x^2}(3-2\sin x)$（$|x|\to\infty$），$\frac{\sin2x}{x\sqrt x}$（$x\to0^+$），$\frac{\ln(1+x)}x$（$x\to0$）。
\end{exercise}
\begin{exercise}{51}{非负数列的极限推断}
设非负数列 $x_n\to0$、$y_n\to1$、$z_n\to+\infty$。证明 $x_n<y_n<z_n$ 最终成立，说明为何不能要求对所有 $n$ 成立；判断 $x_nz_n$ 与 $y_nz_n$ 的变化是否由条件唯一确定，未确定时给出两种不同结果。
\end{exercise}
\begin{exercise}{52}{由组合极限反推函数}
已知 $\lim_{x\to0}\left(1+2x+\frac{f(x)}x\right)=4$。求 $\lim_{x\to0}f(x)$ 及 $\lim_{x\to0}f(x)/x$，并用小量形式写出 $f(x)$；说明此条件能否确定 $f(0)$。
\end{exercise}
\begin{exercise}{53}{无界乘子与反例}
设 $\alpha(x)\to0$（$x\to0$）。证明有界函数 $b(x)$ 满足 $\alpha(x)b(x)\to0$；若仅知 $b$ 在任意去心邻域内无界，分别构造乘积趋零、趋于非零常数和无极限的例子。
\end{exercise}
\begin{exercise}{54}{倒数关系的定义证明}
设 $f$ 在 $x_0$ 的某去心邻域内非零。用定义证明 $|f(x)|\to\infty$ 当且仅当 $1/f(x)\to0$。再构造一个满足上述结论却既不趋于 $+\infty$ 也不趋于 $-\infty$ 的函数。
\end{exercise}
\subsection{极限的运算法则}
\begin{exercise}{55}{无穷大运算的条件}
设 $|f(x)|\to\infty$、$|g(x)|\to\infty$（$x\to x_0$）。判断 $f+g$ 必为无穷大、$f-g\to0$、$kf$ 必为无穷大（$k\ne0$）、$1/(f+g)\to0$ 四个断言，分别证明或举反例；再说明增加 $f,g\to+\infty$ 后哪些结论改变。
\end{exercise}
\begin{exercise}{56}{有限乘积的极限与误差}
求 $\lim_{|x|\to\infty}(1-1/x)(2+2/x)(3-3/x)$；记乘积为 $P(x)$，再求 $\lim_{|x|\to\infty}x(P(x)-6)$，解释两次求极限所需信息的差异。
\end{exercise}
\begin{exercise}{57}{直接代入的合法性}
求 $\lim_{x\to1/2}\frac{2\sin x+2\pi/x}{\cos2x}$，写明分母非零的依据；进一步讨论该表达式在 $x=0$ 处的左右变化趋势。
\end{exercise}
\begin{exercise}{58}{差式极限与有理化}
求 $\lim_{x\to1}\left(\frac1{x+1}-\frac3{x^3+1}\right)$ 与 $\lim_{n\to\infty}(\sqrt{n+2\sqrt n}-\sqrt{n-\sqrt n})$。将第二式中的 $2,-1$ 换成实数 $a,b$，求一般极限。
\end{exercise}
\begin{exercise}{59}{零比零型的消因子}
求 $\lim_{x\to2}\frac{x^2+x-6}{x^2-x-2}$ 与 $\lim_{x\to0}\frac{x^2+x}{\sqrt{1+x}-1}$；分别补定义函数在相应点的值使其连续，并说明约分、乘共轭式时应保留的定义域限制。
\end{exercise}
\begin{exercise}{60}{振荡项与最高次幂}
求 $\lim_{|x|\to\infty}\frac{2x^2+x\cos x}{x^3-x\sin x+1}$，并证明分母在 $|x|$ 充分大时非零；再求原商乘以 $x$ 后的极限。
\end{exercise}
\begin{exercise}{61}{有限和与几何级数}
求 $\lim_{n\to\infty}\sum_{k=1}^n\frac{k}{n^2+1}$ 与 $\lim_{n\to\infty}\sum_{k=1}^na^{-k}$（$a>1$）；再分类讨论第二个极限在 $a\in\mathbb R\setminus\{0\}$ 下存在有限值的全部情形。
\end{exercise}
\begin{exercise}{62}{望远镜求和与差式}
求 $\lim_{n\to\infty}\sum_{k=1}^n\ln\frac{k+1}{k+2}$ 及 $\lim_{n\to\infty}\left(\sqrt{\sum_{k=1}^{n+1}k}-\sqrt{\sum_{k=1}^nk}\right)$，区分有限极限与无穷趋势，并说明能否把第一个式子的极限移入不断增加的各项中。
\end{exercise}
\begin{exercise}{63}{乘积的对数化}
设 $f(x)=2^x$。求 $\lim_{n\to\infty}\frac1{n^2}\ln\prod_{k=1}^nf(k)$；再将 $f(x)$ 改为 $2^{x+c}$，证明该极限与固定常数 $c$ 无关。
\end{exercise}
\begin{exercise}{64}{乘积恒等式与参数范围}
设 $|a|<1$，求 $\lim_{n\to\infty}\prod_{k=0}^n(1+a^{2^k})$。先证明有限乘积恒等式，再讨论 $a=1$、$a=-1$ 两个边界参数，说明不能直接沿用内部参数的化简结论。
\end{exercise}
\begin{exercise}{65}{分式约分后的单侧极限}
讨论 $f(x)=\frac{x^2-1}{x+1}e^{1/(x+1)}$ 在 $x\to-1^-$、$x\to-1^+$ 时的变化，判断双侧极限是否存在；说明消去公因子是否消除了函数在 $-1$ 附近的全部奇异性。
\end{exercise}
\begin{exercise}{66}{根式与无穷远方向}
分别求 $\sqrt{x^2+x+3}-\sqrt{x^2-x+3}$ 在 $x\to+\infty$、$x\to-\infty$ 时的极限，判断 $|x|\to\infty$ 时是否有极限，并指出提取 $\sqrt{x^2}$ 时的注意事项。
\end{exercise}
\begin{exercise}{67}{幂指函数的连续运算}
求 $\lim_{x\to e}(1+\sin2x)^{\ln\ln x}$，验证实数幂在 $e$ 的某邻域内有定义；再把指数改为 $c+\ln\ln x$，求极限关于实数 $c$ 的表达式。
\end{exercise}
\begin{exercise}{68}{底数趋于非单位常数的幂指极限}
求 $\left(\frac{x^2+x\cos x-1}{2x^2-x\sin x}\right)^{x^2+1}$ 与 $\left(\frac{3x+2}{2x-1}\right)^{x^2}$ 在 $|x|\to\infty$ 时的极限或无穷趋势。证明底数最终为正，并用固定常数夹住底数来论证。
\end{exercise}
\begin{exercise}{69}{有限极限反求参数}
已知 $\lim_{x\to2}\frac{x^2+ax+b}{x^2-4}=3$，求 $a,b$。先证明分子在 $2$ 处必须为零，再求参数；并写出使该商在 $2$ 处连续的补定义。
\end{exercise}
\begin{exercise}{70}{对数分母的极限反问题}
设 $a>0$，求全部实数 $a,b$，使 $\lim_{x\to0}\frac{\sin2x}{\ln(a+x)}(\cos x-b)=4$。分别讨论 $a=1$ 与 $a\ne1$，验证所得参数使原商在充分小去心邻域内有定义。
\end{exercise}
\begin{exercise}{71}{无穷远根式抵消的参数}
设 $a>0$ 且 $\lim_{x\to-\infty}(4x+\sqrt{ax^2-bx-1})=1$。求 $a,b$，先比较一次增长项，再用有理化求剩余常数；检查改为 $x\to+\infty$ 时能否有相同结论。
\end{exercise}
\begin{exercise}{72}{根式渐近展开反求系数}
求常数 $a,b$，使 $\lim_{|x|\to\infty}(\sqrt{x^4+2x^2+3}-ax^2-b)=0$；再求所得差式乘以 $x^2$ 后的极限，说明前两个系数确定后尚余何阶误差。
\end{exercise}
\begin{exercise}{73}{运算法则的逆向条件}
设 $f(x)+g(x)\to3$、$f(x)g(x)\to2$（$x\to0$）。证明 $f,g$ 在某去心邻域内有界；构造两者均无极限的例子。若再知该邻域内 $f(x)\le g(x)$，求两者极限。
\end{exercise}
\begin{exercise}{74}{根式抵消的参数分类}
对实数 $a,b$，分类讨论 $\lim_{x\to+\infty}(\sqrt{x^2+ax+b}-x)$，并求 $\lim_{x\to+\infty}x(\sqrt{x^2+ax+b}-x-a/2)$。证明表达式在充分大的 $x$ 上有实数意义。
\end{exercise}
\subsection{极限存在准则与两个重要极限}
\begin{exercise}{75}{阶乘增长与夹逼}
求 $\lim_{n\to\infty}\frac{3^n}{n!}\sin\frac1n$；不用斯特林公式，证明存在常数 $C>0$，使充分大的 $n$ 满足 $3^n/n!\le C(1/2)^n$。
\end{exercise}
\begin{exercise}{76}{变上限和的统一估计}
证明 $\lim_{n\to\infty}\sum_{k=1}^n\frac{k}{n^2+2n+k}=\frac12$；给出上下界，并由此估计该和与 $1/2$ 的误差为 $O(1/n)$。
\end{exercise}
\begin{exercise}{77}{逐渐缩小相位差的夹逼}
求 $\lim_{x\to+\infty}\bigl(\cos\sqrt{x+1}-\cos\sqrt x\bigr)$。用三角恒等式和根式差估计证明，不得分别给两个振荡项取极限；再讨论将 $1$ 换成固定实数 $a$ 的情形。
\end{exercise}
\begin{exercise}{78}{递推数列的收敛证明}
设 $x_1>0$、$x_{n+1}=\frac12(x_n+4/x_n)$。证明从第二项起 $x_n\ge2$ 且单调，求其极限；再导出 $x_{n+1}-2$ 与 $x_n-2$ 的精确关系，说明初值小于 $2$ 时第一步的变化。
\end{exercise}
\begin{exercise}{79}{指数递推与不等式}
设 $x_1>0$、$x_ne^{x_{n+1}}=e^{x_n}-1$。证明递推式始终确定正数 $x_{n+1}$，且数列严格递减、收敛；求其极限。可使用 $e^t>1+t$（$t\ne0$），不得先假设极限存在。
\end{exercise}
\begin{exercise}{80}{积分数列的递推与比值极限}
设 $a_n=\int_0^1x^n\sqrt{1-x^2}\,\dd x$（$n\ge0$）。证明 $a_n$ 严格递减且 $a_n=\frac{n-1}{n+2}a_{n-2}$（$n\ge2$）；利用递推式和单调性夹逼求 $\lim_{n\to\infty}a_n/a_{n-1}$，写明正性保证的除法条件。
\end{exercise}
\begin{exercise}{81}{重要极限与无穷远估计}
设 $f(x)=\sin2x/(5x)$（$x\ne0$）。分别求 $x\to0$ 与 $|x|\to\infty$ 时的极限，并给出所用依据；说明为何同一分式在两个变化过程中需要不同的估计。
\end{exercise}
\begin{exercise}{82}{重要极限与参数}
求 $\lim_{x\to0}\frac{\tan3x}{4x}$；推广求 $\lim_{x\to0}\frac{\tan(ax)}{bx}$（$a\in\mathbb R,b\ne0$），并说明 $a=0$ 时结论是否仍适用。
\end{exercise}
\begin{exercise}{83}{反三角函数的换元极限}
不用洛必达法则，求 $\lim_{x\to0}\frac{\arcsin2x}{x}$ 和 $\lim_{x\to0}\frac{\arctan x}{2x}$；分别说明换元后自变量趋于零的依据及反三角函数的取值范围。
\end{exercise}
\begin{exercise}{84}{同一表达式的不同变化过程}
分别求 $\frac{2x-\sin3x}{2x+\sin4x}$ 在 $x\to0$ 与 $|x|\to\infty$ 时的极限。再确定实数 $a$ 的全部取值，使 $\lim_{x\to0}\frac{ax-\sin3x}{ax+\sin4x}$ 有有限值，并求该值。
\end{exercise}
\begin{exercise}{85}{复合零点的因式分解}
求 $\lim_{x\to3}\frac{\sin(x^2-2x-3)}{x-3}$；再讨论 $x\to-1$ 时同一分式的极限。将正弦的自变量因式分解，说明两个零点处的分母是否同时趋零。
\end{exercise}
\begin{exercise}{86}{余弦差与阶数}
求 $\lim_{x\to0}\frac{1-\cos2x}{x^2}$；进一步分类讨论 $\lim_{x\to0}\frac{1-\cos2x}{|x|^p}$（$p\in\mathbb R$），说明何时为非零有限值。
\end{exercise}
\begin{exercise}{87}{固定参数与变动参数}
对固定实数 $x$，求 $\lim_{n\to\infty}n^2\sin\frac{x^2+1}{n^2}$。再令 $x=n$，判断能否把固定参数的结论直接用于此时的表达式，并求该表达式除以 $n^2$ 后的极限。
\end{exercise}
\begin{exercise}{88}{幂指型极限}
求 $\lim_{|x|\to\infty}(1-2/x)^{2x+100}$，证明底数最终为正；再将指数中的 $100$ 换为 $c x$，求极限关于实数 $c$ 的表达式。
\end{exercise}
\begin{exercise}{89}{商底数的幂指极限}
求 $\lim_{|x|\to\infty}\left(\frac{x+1}{x-1}\right)^{2x}$；推广求 $\lim_{|x|\to\infty}\left(\frac{x+a}{x+b}\right)^{cx}$，其中 $a,b,c$ 为固定实数，并说明实数幂的定义条件最终成立。
\end{exercise}
\begin{exercise}{90}{二次尺度的重要极限}
求 $\lim_{|x|\to\infty}\left(\frac{2x^2-1}{2x^2+1}\right)^{x^2}$。若把指数 $x^2$ 改为 $|x|^p$，对实数 $p$ 分类讨论极限。
\end{exercise}
\begin{exercise}{91}{对数等价量与参数}
设 $k\ne0$，求 $\lim_{x\to0}\frac{\sin3x}{\ln(1+kx)}$；证明对数在 $0$ 的某邻域内有定义。再讨论当 $k>0$、$x\to+\infty$ 时该商的极限，区分局部等价量的适用过程。
\end{exercise}
\begin{exercise}{92}{取对数与符号控制}
求 $\lim_{x\to0}(1+\tan2x)^{-3/x}$，并验证底数在 $0$ 的某邻域内为正；推广求 $(1+\tan ax)^{b/x}$ 的极限，其中 $a,b\in\mathbb R$。
\end{exercise}
\begin{exercise}{93}{指数与底数的共同变化}
求 $\lim_{x\to0}(\cos x)^{-1/\sin x}$ 和 $\lim_{x\to0}\left(\frac{\tan x}{x}\right)^{1/x}$。将两式的指数分别改为 $-1/\sin^2x$ 和 $1/x^2$，再求极限，比较一阶与二阶尺度。
\end{exercise}
\begin{exercise}{94}{幂平均的局部极限}
设 $a,b>0$，求 $\lim_{x\to0}\left(\frac{a^x+b^x}{2}\right)^{1/x}$；推广到 $\left(\lambda a^x+(1-\lambda)b^x\right)^{1/x}$（$0\le\lambda\le1$），说明端点权重情形。
\end{exercise}
\begin{exercise}{95}{对数与幂的括号范围}
设 $a\ne\frac12$，求 $\lim_{n\to\infty}\ln\left[\left(\frac{n-2an-1}{n(1-2a)}\right)^n\right]$，说明对数为何最终有定义；再求幂本身的极限。
\end{exercise}
\begin{exercise}{96}{幂指极限反求参数}
求全部实数对 $(a,b)$，使 $\lim_{|x|\to\infty}\left(\frac{x^2}{(x-a)(x+b)}\right)^x=e^2$。证明所得条件充分且必要，并检查底数最终为正。
\end{exercise}
\begin{exercise}{97}{递推极限与初值}
设 $x_1\ge0$、$x_{n+1}=\sqrt{2+x_n}$。分别讨论 $0\le x_1<2$、$x_1=2$、$x_1>2$ 的单调性，证明数列收敛并求极限；当 $x_1\ne2$ 时给出 $|x_{n+1}-2|/|x_n-2|$ 的极限。
\end{exercise}
\begin{exercise}{98}{重要极限中的高阶修正}
对实数 $a,b$，求 $\lim_{x\to0}(1+ax+bx^2)^{1/x}$；当 $a=0$ 时，再求 $\lim_{x\to0}(1+bx^2)^{1/x^2}$。各情形均须说明底数在充分小邻域内的正性。
\end{exercise}
\subsection{无穷小的比较}
\begin{exercise}{99}{以根式为基准比较阶数}
在 $x\to0^+$ 时，分别求 $e^{\sqrt x}-1$、$\ln(1+\sqrt x)$、$1-\cos\sqrt x$、$\arctan\sqrt x$ 与 $\sqrt x$ 的比值极限；对不等价者给出其关于 $\sqrt x$ 的阶数与首项。
\end{exercise}
\begin{exercise}{100}{不同阶无穷小的比较}
当 $x\to0$ 时，分别求 $\arctan2x$、$\sqrt{1+x}-1$、$1-\cos2x$、$\tan x-\sin x$ 的首个非零等价项，按阶数排序并指出最高阶者；说明同阶是否必然等价。
\end{exercise}
\begin{exercise}{101}{无穷远的等价量}
求 $\lim_{|x|\to\infty}\frac{2x-3}{x^2-x+1}$；进一步求 $x\frac{2x-3}{x^2-x+1}$ 的极限，并写出原商相对于 $1/x$ 的等价关系。
\end{exercise}
\begin{exercise}{102}{指数差的消去阶数}
当 $x\to0$ 时，确定使 $e^{x\cos x^2}-e^x$ 与 $x^n$ 同阶的正整数 $n$，并求两者比值的极限；论证时须保留指数差中首先不为零的项。
\end{exercise}
\begin{exercise}{103}{三阶等价无穷小}
证明 $\tan x-x\sim x^3/3$（$x\to0$）；利用此结论求 $\lim_{x\to0}\frac{\tan(ax)-a\tan x}{x^3}$，并讨论实数 $a$ 使极限为零的全部情形。
\end{exercise}
\begin{exercise}{104}{高阶无穷小与首项}
证明 $(2^{x^2}-1)\arctan2x=o(x^2)$（$x\to0$）；再求其与 $x^3$ 的比值极限，说明“高于二阶”是否意味着“高于三阶”。
\end{exercise}
\begin{exercise}{105}{等价关系能确定的参数信息}
求全部实数 $a,b,c$，使 $\frac1{ax^2+bx+c}\sim\frac1{2x^2+1}$（$|x|\to\infty$）。若进一步要求两式之差为 $o(x^{-4})$，再求参数；说明增强前后分别约束了哪些系数。
\end{exercise}
\begin{exercise}{106}{含实参数幂的阶数比较}
当 $x\to0^+$ 时，求全部实数 $k$，使 $(\sqrt{1+x^2}-1)\ln(1+x^2)=o\bigl(x(\arcsin x)^k\bigr)$；再讨论临界参数时的比值极限。
\end{exercise}
\begin{exercise}{107}{乘积中的高阶消去}
求 $\lim_{x\to0}\frac{(\sin x-\tan x)(e-e^{\cos x})}{(2x+1)\sin^5x}$。分别写出分子两个因子的等价首项，并说明为何不能把 $\sin x$ 与 $\tan x$ 都代换成 $x$ 后相减。
\end{exercise}
\begin{exercise}{108}{等价关系反求两个参数}
设 $ab\ne0$，当 $x\to0$ 时 $x-\sin(ax)\sim x^2\ln(1-bx)$。求 $a,b$ 并验证；说明为何必须先排除分子的一阶项，再比较三阶项。
\end{exercise}
\begin{exercise}{109}{对数极限反推原函数}
设 $a>0$、$a\ne1$，且 $\lim_{x\to0}\frac{\ln(1+f(x)/\sin^2x)}{a^2-1}=A\in\mathbb R$，原式在充分小去心邻域内有定义。求 $\lim_{x\to0}f(x)/x^2$；分别讨论 $A=0$ 与 $A\ne0$ 时能否直接使用 $\ln(1+u)\sim u$。
\end{exercise}
\begin{exercise}{110}{已知高阶余项反求极限}
已知 $\lim_{x\to0}\frac{\sin6x+xf(x)}{x^3}=0$，求 $\lim_{x\to0}\frac{6+f(x)}{x^2}$；再写出 $f(x)$ 到 $x^2$ 阶的渐近表达式，并说明是否能据此确定 $f(0)$。
\end{exercise}
\begin{exercise}{111}{振荡余项的阶估计}
求 $\lim_{x\to0}\frac{4\sin x+3x^2\cos(1/x^2)}{(1+\cos x)\ln(1+2x)}$。若把分子中的 $x^2$ 改为 $|x|^p$，确定极限仍等于原值的全部实数 $p$，对边界情形说明理由。
\end{exercise}
\begin{exercise}{112}{根式差与三阶比较}
求 $\lim_{x\to0}\frac{\sqrt{1+\tan x}-\sqrt{1+\sin x}}{x\ln(1+x)-x^2}$；指出两个差式各自的首个非零阶，并验证根式和对数在充分小邻域内有定义。
\end{exercise}
\begin{exercise}{113}{指数相近时的四阶消去}
求 $\lim_{x\to0}\frac{e^{x^2}-e^{2-2\cos x}}{x^3\sin x}$；写出两个指数之差的首个非零项，并说明分别只保留指数到二阶会丢失什么信息。
\end{exercise}
\begin{exercise}{114}{对数差的等价转化}
求 $\lim_{x\to0}\frac{\ln(\sin^2x+e^{2x})-2x}{\ln(x^2+e^{3x})-3x}$。先把每个对数差改写为 $\ln(1+u)$ 型，再说明相减结构下允许等价代换的具体位置。
\end{exercise}
\begin{exercise}{115}{无穷远复合等价量}
求 $\lim_{|x|\to\infty}x\left[\sin\ln(1+4/x)-\sin\ln(1+1/x)\right]$；将 $4,1$ 换为实数 $a,b$，求一般结果，并证明两个对数最终有定义。
\end{exercise}
\begin{exercise}{116}{积分无穷小的阶数比较}
在 $x\to0^+$ 时，求 $\int_0^x(e^t-1)\,\dd t$、$\int_0^x\ln(1+\sqrt{t^2})\,\dd t$、$\int_0^x\sin(t^2)\,\dd t$、$\int_0^x\sqrt{\sin^2t}\,\dd t$ 的等价首项，按阶数排序并找出最高阶者。说明两个根式在正半轴上如何化简。
\end{exercise}
\begin{exercise}{117}{极限减去常数后的阶数}
设 $f(x)=\frac{1+x}{\sin x}-\frac1x$，$a=\lim_{x\to0}f(x)$。求 $a$；确定使 $f(x)-a$ 与 $x^k$ 同阶的正整数 $k$，并求比值极限；说明分子分母主项相消后应保留到哪一阶。
\end{exercise}
\begin{exercise}{118}{积分分母的非零有限极限}
求全部常数 $a,b,c$（$c\ne0$），使 $\lim_{x\to0}\frac{ax-\sin(bx)}{\int_0^x\ln(1+t^2)\,\dd t}=c$。先比较阶数再确定系数关系，并验证所求关系足以保证极限存在。
\end{exercise}
\begin{exercise}{119}{加减式中等价代换的限度}
已知 $\alpha\sim\beta$、$\gamma\sim\delta$（$x\to0$），说明是否必有 $\alpha-\gamma\sim\beta-\delta$，并给出反例。再求 $\lim_{x\to0}\frac{\sin x-x\cos x}{\tan x-x}$，用首项说明相消时需要保留到哪一阶。
\end{exercise}
\begin{exercise}{120}{参数控制的阶数跃迁}
设 $F_a(x)=\ln(1+x)-x+ax^2$。对所有实数 $a$，求 $F_a(x)$ 在 $x\to0$ 时的首个非零等价项，并确定其阶数；求全部 $(a,c)$，使 $F_a(x)\sim cx^3$ 且 $c\ne0$。
\end{exercise}
\subsection{函数的连续性与间断点}
\begin{exercise}{121}{差商极限与连续性}
已知 $f$ 在 $1$ 处连续，且 $\lim_{x\to1}\frac{f(x)-4}{x-1}=A\in\mathbb R$。求 $f(1)$，证明 $f$ 在 $1$ 处可导并求其导数；去掉连续性条件后，说明极限条件还可以推出哪些结论。
\end{exercise}
\begin{exercise}{122}{连续性的定义证明}
用 $\varepsilon$--$\delta$ 定义证明 $f(x)=2x^2-x+1$ 在任意实数 $x_0$ 处连续，并给出依赖 $\varepsilon,x_0$ 的明确 $\delta$；说明该估计是否可在固定闭区间上统一选取。
\end{exercise}
\begin{exercise}{123}{连续延拓及原式的定义范围}
设 $f(x)=\begin{cases}(e^{2x}-1)/\sin4x,&x\ne0,\\k,&x=0.\end{cases}$ 在 $0$ 的充分小邻域内定义。求使其在 $0$ 处连续的 $k$；列出非零分母零点，判断仅确定 $k$ 能否使原表达式在整个 $\mathbb R$ 上有定义且连续。
\end{exercise}
\begin{exercise}{124}{分段点的连续性与单侧修复}
设 $f(x)=\begin{cases}\frac{\sin2x}{\tan x}(1+x^2),&x>0,\\2,&x=0,\\\ln\sqrt{1-2x},&x<0.\end{cases}$，各支在 $0$ 的充分小邻域内考虑。判断 $f$ 在 $0$ 处的连续性；若只允许给左支加常数 $c$，求能使函数连续的全部 $c$。
\end{exercise}
\begin{exercise}{125}{振荡因子的连续延拓}
讨论 $f(x)=\begin{cases}(x-1)^2\cos\frac1{x-1},&x<1,\\(x+1)\ln x,&x\ge1\end{cases}$ 在定义域内的连续性；若把左支中的 $(x-1)^2$ 改为 $|x-1|^p$，确定在 $1$ 处连续的全部实数 $p$。
\end{exercise}
\begin{exercise}{126}{根式与反三角函数的延拓}
为 $f(x)=\frac{\sqrt{1+x^2}-\sqrt{1-x^2}}{x\arcsin x}$ 补定义 $f(0)$，使其在 $0$ 处连续；写出延拓后函数的定义域，并判断它在定义域的端点是否连续。
\end{exercise}
\begin{exercise}{127}{由数列极限定义的函数}
对 $x>0$，设 $f(x)=\lim_{n\to\infty}\frac1n\ln(e^{nx}+x^2)$。求 $f(x)$ 并讨论其连续性；若按同一极限式将定义范围扩大到 $\mathbb R$，求扩展后的分段表达式并判断 $0$ 处是否连续。
\end{exercise}
\begin{exercise}{128}{连续函数值的反推}
设 $f$ 在 $0$ 的某邻域内连续，在充分小的 $x>0$ 时原商有定义，且 $\lim_{x\to0^+}\frac{1-\cos(\sqrt x\,f(x))}{(e^{2x}-1)f(x)}=\frac12$。求 $f(0)$；证明时须同时考虑 $f(0)=0$ 与 $f(0)\ne0$，不能未经论证就约去可能趋零的量。
\end{exercise}
\begin{exercise}{129}{有理函数的间断点}
求 $f(x)=\frac{x^2-1}{x^2-3x+2}$ 的全部间断点并分类；对可去间断点作连续延拓，对其余间断点写出左右变化趋势，说明二者不能采用同一补定义方式处理。
\end{exercise}
\begin{exercise}{130}{指数分式的间断点分类}
求 $f(x)=\frac1{e^{x/(x-1)}-1}$ 的全部间断点，计算各自的左右极限或无穷趋势并分类；判断每一点能否通过补定义成为连续点。
\end{exercise}
\begin{exercise}{131}{可去间断点的参数条件}
设 $f(x)=\begin{cases}\frac{x^2\sin(1/(2x))}{2^x-1},&x<0,\\b,&x=0,\\a+\frac{\ln(1+3x)}x,&x>0.\end{cases}$ 求使 $0$ 为可去间断点的全部参数对 $(a,b)$；再求使其在 $0$ 连续的参数对，说明可去间断点与连续点对 $b$ 的要求。
\end{exercise}
\begin{exercise}{132}{点态极限与间断点}
设 $f(x)=\lim_{n\to\infty}\frac{1+x}{1+x^{2n}}$。求其分段表达式、全部连续区间及间断点类型；说明为何不能从各个近似函数连续直接推断 $f$ 连续。
\end{exercise}
\begin{exercise}{133}{幂平均极限与最大值函数}
设 $a_1,\ldots,a_m>0$，证明 $\lim_{n\to\infty}(a_1^n+\cdots+a_m^n)^{1/n}=\max_{1\le i\le m}a_i$。据此求 $f(x)=\lim_{n\to\infty}\bigl(1+x^n+(x^2/2)^n\bigr)^{1/n}$（$x\ge0$）的分段表达式，并讨论其连续性。
\end{exercise}
\begin{exercise}{134}{连续与可去间断的参数分类}
设 $f(x)=\begin{cases}\frac{\ln(1+ax^3)}{x-\arcsin x},&x<0,\\6,&x=0,\\\frac{e^{ax}+x^2-ax-1}{x\tan(x/4)},&x>0,\end{cases}$ 各支在 $0$ 的充分小邻域内考虑。求使 $f$ 在 $0$ 连续的全部 $a$，以及使 $0$ 为可去间断点的全部 $a$；给出各情形的左右极限。
\end{exercise}
\begin{exercise}{135}{参数极限生成函数}
对 $\sin x\ne0$ 的实数 $x$，设 $f(x)=\lim_{t\to x}\left(\frac{\sin t}{\sin x}\right)^{x/(\sin t-\sin x)}$，极限沿表达式有定义的 $t$ 取。求 $f(x)$，找出其定义域边界上的间断点并分类；区分 $x=0$ 与其他 $\pi$ 的整数倍。
\end{exercise}
\begin{exercise}{136}{稠密分段与连续点}
设 $f(x)=\begin{cases}x^2,&x\in\mathbb Q,\\ax+b,&x\notin\mathbb Q.\end{cases}$ 求其全部连续点关于参数 $a,b$ 的分类；给出恰有两个连续点的参数，并证明该函数不可能在任何非退化区间上连续。
\end{exercise}
\begin{exercise}{137}{振荡函数的参数延拓}
设 $f_{p,c}(x)=\begin{cases}|x|^p\sin(1/x),&x\ne0,\\c,&x=0.\end{cases}$ 对全部实数 $p,c$，讨论 $0$ 处的连续性与间断点类型；对 $p<0$，说明为什么函数在每个邻域内无界仍不构成无穷间断点的充分条件。
\end{exercise}
\subsection{连续函数的运算与初等函数的连续性}
\begin{exercise}{138}{连续运算的分母条件}
设 $f(x)=\frac{\lg(100+x^2)}{a^2+\arcsin2x}$。分别讨论 $a\ne0$ 与 $a=0$ 时 $x\to0$ 的双侧极限或单侧趋势，说明何时可直接使用连续函数商的法则。
\end{exercise}
\begin{exercise}{139}{间断性在运算中的保留}
设 $f$ 在 $\mathbb R$ 上连续且处处非零，$\varphi:\mathbb R\to\mathbb R$ 有间断点。判断 $\varphi\circ f$、$f\circ\varphi$、$\varphi^2$、$\varphi/f$ 是否必有间断点；必然成立者证明，其余给出满足全部条件的反例。
\end{exercise}
\begin{exercise}{140}{反函数连续性的条件}
设 $f$ 在闭区间 $[a,b]$ 上连续且严格递增。证明 $f^{-1}$ 在 $[f(a),f(b)]$ 上连续，端点按单侧连续理解；举例说明仅有严格递增而无连续性时，函数值域可能不是区间。
\end{exercise}
\begin{exercise}{141}{连续复合与零点集合}
设 $f$ 在 $[a,b]$ 上连续。证明 $g(x)=\max\{f(x),0\}$ 与 $h(x)=|f(x)|$ 连续；再证明若 $f$ 处处非零，则 $1/f$ 在该区间上连续且有界，并指出闭区间条件在有界性结论中的作用。
\end{exercise}
\subsection{闭区间上连续函数的性质}
\begin{exercise}{142}{三次方程的实根个数}
证明方程 $x^3-9x-2=0$ 恰有三个不同实根，并将每个根定位在一个长度为 $1$ 且互不相交的开区间内；分别说明“至少三个”与“至多三个”的依据。
\end{exercise}
\begin{exercise}{143}{零点存在与唯一性}
证明方程 $x=\frac12\cos x$ 在 $\mathbb R$ 上恰有一个实根，并给出长度不超过 $1/2$ 的含根闭区间；存在性与唯一性须分别论证。
\end{exercise}
\begin{exercise}{144}{平移差函数与零点}
设 $f$ 在 $[0,1]$ 上连续且 $f(0)=f(1)$。证明存在 $\xi\in[0,1/2]$，使 $f(\xi+1/2)=f(\xi)$；说明为何平移差函数只能先定义在 $[0,1/2]$ 上。
\end{exercise}
\begin{exercise}{145}{半无限区间的有界性}
设 $f$ 在 $[a,+\infty)$ 上连续，且 $\lim_{x\to+\infty}f(x)=L\in\mathbb R$。证明 $f$ 有界；再分别举例说明它未必取得最大值、未必取得最小值。
\end{exercise}
\begin{exercise}{146}{加权平均与介值性}
设 $f$ 在 $[a,b]$ 上连续，$a<x_1<x_2<x_3<b$。证明存在 $\xi\in(a,b)$ 使 $f(x_1)+2f(x_2)+3f(x_3)=6f(\xi)$；若三点函数值不全相同，证明可取 $\xi\in(x_1,x_3)$。
\end{exercise}
\begin{exercise}{147}{开区间端点与有界性}
设 $f(x)=\frac{x\sin(x-2)}{|x|(x-1)(x-2)^2}$。分别判断它在 $(-1,0)$、$(0,1)$、$(1,2)$、$(2,3)$ 内是否有界；有界情形用连续延拓或明确估计证明，无界情形给出导致无界的端点趋势。
\end{exercise}
\begin{exercise}{148}{等分平移与望远镜求和}
设 $f$ 在 $[0,1]$ 上连续且 $f(0)=f(1)$。证明对每个正整数 $n$，存在 $x_0\in[0,1-1/n]$ 使 $f(x_0)=f(x_0+1/n)$。可考虑平移差函数在等分点上的和，并说明 $n=1$ 的情形。
\end{exercise}
\begin{exercise}{149}{闭区间自映射的不动点}
设连续函数 $f:[a,b]\to[a,b]$，其中 $a<b$。证明至少存在一点 $\xi\in[a,b]$ 使 $f(\xi)=\xi$；若另有 $|f(x)-f(y)|\le q|x-y|$（$0\le q<1$），证明不动点唯一；举例说明仅连续时唯一性可能不成立。
\end{exercise}
\begin{exercise}{150}{严格不等式的统一余量}
设 $f,g$ 在 $[a,b]$ 上连续且处处 $f(x)<g(x)$。证明存在 $\eta>0$，使所有 $x\in[a,b]$ 均满足 $f(x)+\eta\le g(x)$；再给出开区间 $(a,b)$ 上的反例，说明逐点严格不等式未必给出统一正余量。
\end{exercise}
